\subsection{Meeting of 27.01.26}
\subsubsection{Since Last Time}
\begin{itemize}
    \item Read through papers
    \item Understand diffusion
\end{itemize}
\subsubsection{Questions}
\begin{itemize}
    \item Process:
    \begin{itemize}
        \item Bi-weekly meetings? \comment{Approx. every 2 weeks; via Zoom or in person}
        \item Accompanying seminar? \comment{There is this seminar, 3 sessions; write proposal; registrations completely independent of this; coordinate dates with GR so he can give feedback too + send proposal to the group}
        \item Private repo/code environment from the group? Or generally shared resources somewhere, or should I keep everything on my end? \comment{GitHub repository}
    \end{itemize}
    \item Content-wise
    \begin{itemize}
        \item Focus on destroy sets, correct? Is repair easier than destroy, or why destroy? \comment{New implementation is coming anyway, }
        \item Meaningful to look at current codebase for the DIFUSCO paper again, or not? It was mentioned that there were limitations in the python impl., those would interest me. \comment{Not meaningful at the moment, implementation should rather start from scratch}
        \item Question for myself: Regarding loss function/metric for good destroy sets: The unpublished paper DIFUSCO-LNS trains a set of predefined destroy sets, which were found using Local Branching. In the context of DynEADARP this looks a bit different of course. That is naturally the core question. Any ideas/approaches here? Do you think about something like this before implementation or only at the end?
    \end{itemize}
\end{itemize}
\subsubsection{General Notes}
\begin{itemize}
    \item Julia implementation of diffusion approach is definitely the first goal. First problem setting not important; possibly MISP first. There is already a bit of codebase for this.
    \item Formal template for the thesis is available from the dean's office
    \item There are GPUs from the group (GPU01@actuwien.ac.at;CONAN@actuwien.ac.at)
    \item CPU Cluster in the AC was also already used for HeuOpt
    \item Place working paper in the repository so everyone has access
    \item Use VSC for Julia with plugin
    \item Look at FLUX (is a NN package) + Graphneuralnetwork.jl also a good package to get to know; important for the implementation.
    \item Important trick: Positional embedding; difusco confuses here.
    \item Unit tests for individual components are probably also meaningful.
    \item Fundamental problem with the DIFUSCO implementation is the copying around of data - would have to analyze with profiler; Question: Profiler?
    \item General process: First make it work and only then look at optimization and GPU usage
    \item Building on DIFUSCO there are a few other papers: Question: How many denoising steps do I need. That can be optimized, GR is sending me paper.
    \item Tackle only one DIFUSCO approach (discrete noise prob. meant) afterwards one can still optimize other things.
\end{itemize}